\section{Background}
The \texttt{postwfn} module performs calculations that require electronic-state wave functions rather than MOs alone. A \texttt{\$\$statepack} block supplies the orbitals and the selected electronic-state data, as described in Section~\ref{sec-statepack}. The selected states may be adiabatic states from an external quantum chemistry calculation or quasi-diabatic states previously exported by \progname{}. The main calculations are based on the one-particle density matrix, the transition density matrix, and the difference density matrix.

\paragraph{One-particle density matrices.}
For an electronic state, the one-particle density matrix is defined from the expectation values of one-electron excitation operators and describes the distribution of electrons in the chosen orbital representation. Diagonalization of this matrix gives the natural orbitals and their occupation numbers. By contracting the density matrix with the corresponding one-electron orbital integrals, \texttt{postwfn} can calculate permanent one-electron properties, including the electric dipole moment and other supported electric multipole moments. The module can also generate real-space electron- and spin-density data for the selected state and write them to cube files.

\paragraph{Transition density matrices.}
For a pair of electronic states, the transition density matrix is defined from matrix elements of one-electron excitation operators between the two states. Singular-value decomposition of this matrix gives the natural transition orbitals (NTOs) and their associated singular values. By contracting the transition density matrix with the corresponding one-electron orbital integrals, \texttt{postwfn} can calculate interstate one-electron properties, including transition electric dipole moments and other supported transition multipole moments. The module can also generate transition-density data in real space, together with the corresponding particle and hole densities, and write these quantities to cube files.

\paragraph{Difference density matrices.}
For a target state and a reference state, the difference density matrix is defined as the one-particle density matrix of the target state minus that of the reference state. Diagonalization of the difference density matrix gives the natural difference orbitals and their eigenvalues.
%The positive and negative components of the matrix can further be used to construct attachment and detachment densities.
%By contracting the difference density matrix with the corresponding one-electron orbital integrals, differences in supported one-electron properties between the target and reference states can be calculated.
The module can also generate real-space difference, attachment, and detachment densities and write them to cube files. When a specific loaded state is required as the reference, it can be selected with the corresponding \keyref{key-postwfn-refid}{refid} keyword.
